\section{Theoretical Setup}\label{sec:setup}

\subsection{Problem Setting}
Let $(X,\Sigma)$ be a nonempty measurable space and let
$C\subseteq\{0,1\}^{X}$ be a nonempty binary concept class.  Define
\[
x\equiv_C x'\quad\Longleftrightarrow\quad
c(x)=c(x')\ \text{for every }c\in C,
\]
let $Q_C=X/{\equiv_C}$, and write $\kappa:X\to Q_C$ for the quotient map.
The quotient carries the discrete sigma-field $2^{Q_C}$.  For $c\in C$,
write $\bar c(\kappa(x))=c(x)$ and set
$\bar C=\{\bar c:c\in C\}$.  The released hypothesis space and decoder are
\[
 H_C=\{0,1\}^{Q_C},\qquad
 \mathcal H_C=\bigotimes_{q\in Q_C}2^{\{0,1\}},\qquad
 \operatorname{Dec}_C(\bar h)=\bar h\circ\kappa .
\]
The released object is allowed to be improper: $\bar h$ need not lie in
$\bar C$.

Write
\[
 v=\operatorname{VC}(C),\qquad d=\operatorname{LD}(C),
\]
and use natural logarithms.  For $u\ge0$, define $\log^*u=0$ for
$u\le1$ and, for $u>1$, define
$\log^*u=\min\{j\ge1:\log^{\circ j}(u)\le1\}$.  For cardinality use
\[
\log^+|C|=\begin{cases}\max\{1,\log|C|\},&|C|<\infty,\\
 +\infty,&|C|=\infty.
 \end{cases}
\]
Put $Z_X=X\times\{0,1\}$ and $Z_Q=Q_C\times\{0,1\}$, with the product
sigma-fields.  For $N\in\mathbb N_0$, define
\[
 T_N((x_r,y_r)_{r=1}^{N})=((\kappa(x_r),y_r))_{r=1}^{N}.
\]
Raw labeled datasets use replace-one adjacency.  If $D$ is a probability
measure on $(X,\Sigma)$ and $c\in C$, let
\[
 \bar D=\kappa_\#D,\qquad
 P_{D,c}=\mathcal L(x,c(x)),\qquad
 P_{\bar D,\bar c}=\mathcal L(q,\bar c(q)),
\]
where $x\sim D$ and $q\sim\bar D$.  For $\bar h\in H_C$, define
\[
 \operatorname{err}_D(\operatorname{Dec}_C(\bar h),c)
 =D\{x:\bar h(\kappa(x))\ne c(x)\},\qquad
 \operatorname{err}_{\bar D}(\bar h,\bar c)
 =\bar D\{q:\bar h(q)\ne\bar c(q)\}.
\]

A learner with sample size $N$ is a Markov kernel
$A_N:Z_X^N\leadsto H_C$.  It is $(\varepsilon,\delta)$-differentially
private when, for every neighboring $s\sim s'$ and every
$E\in\mathcal H_C$,
\[
 A_N(s,E)\le e^{\varepsilon}A_N(s',E)+\delta .
\]
The conditional-scope realizable private sample complexity is
\[
\begin{aligned}
m_C(\alpha,\beta;\varepsilon,\delta):=\inf\bigl\{N\in\mathbb N_0:\;&\text{there exists a Markov kernel }A_N:
 Z_X^N\leadsto H_C\text{ that is }(\varepsilon,\delta)\text{-DP, and}\\[-2pt]
&\sup_D\sup_{c\in C}\Pr_{S\sim P_{D,c}^{N},\,\bar H\sim A_N(S,\cdot)}
 [\operatorname{err}_D(\operatorname{Dec}_C(\bar H),c)>\alpha]\le\beta\bigr\}.
\end{aligned}
\]
The first supremum ranges over all probability measures on $(X,\Sigma)$;
the DP predicate quantifies over every raw replace-one pair and every
$E\in\mathcal H_C$ as above.  Learners may be computationally unbounded and
need not be proper.

\subsection{Primitive Assumptions}
\begin{assumption}[Finite Littlestone dimension]\label{assump:finite-littlestone}
$C$ is nonempty and $d=\operatorname{LD}(C)<\infty$.  No finite-cardinality,
finite-domain, product, properness, efficiency, or computational assumption
is imposed.
\end{assumption}

\begin{assumption}[Finite-or-countable measurable evaluation quotient]\label{assump:countable-evaluation-quotient}
$Q_C=X/{\equiv_C}$ is finite or countably infinite, and every cell
$\kappa^{-1}(\{q\})$ belongs to $\Sigma$.  Equivalently,
$\kappa:(X,\Sigma)\to(Q_C,2^{Q_C})$ is measurable.  This is a static
pre-sampling condition and does not assume measurability of restrictions,
lists, selectors, events, or learner kernels.
\end{assumption}

\begin{assumption}[Realizable iid sampling]\label{assump:realizable-iid}
For utility, $D$ is arbitrary on $(X,\Sigma)$, $c$ is arbitrary in $C$,
and the labeled sample is iid from $P_{D,c}$.  Privacy is required on every
raw neighboring labeled input, including nonrealizable inputs.
\end{assumption}

\begin{assumption}[Learning, confidence, and approximate-DP parameters]\label{assump:approximate-dp-regime}
$0<\alpha,\beta<1/4$, $0<\varepsilon\le1$, and $0<\delta<1$.  For the
positive-dimensional VC-sensitive comparison, the source convention
$\delta\ll1/N$ means that along the stated asymptotic sequences,
\[
 \delta K\Lambda(d,v,\alpha,\beta,\varepsilon,\delta)^q
 R_{\mathrm{VC}}(d,v,\alpha,\beta,\varepsilon,\delta)\longrightarrow0
\]
for the universal constants in the theorem below.  No fixed positive
$\delta$ is silently absorbed into that source-facing comparison.
\end{assumption}

\subsection{Formalized Goal}
The theorem is explicitly conditional on the four assumptions above.  It
constructs the quotient-first totalized law described in the next section,
pulls it back through $T_N$, and proves measurability, all-input
$(\varepsilon,\delta)$-privacy, and realizable population error at most
$\alpha$ with failure probability at most $\beta$.  The finite-or-countable
quotient is part of the theorem scope; classes with uncountably many
evaluation types remain outside the claim.

Define the three source-certified rate expressions
\[
 R_{\mathrm{VC}}=
 \frac{d^4\bigl(v+\log(1/\beta)\bigr)\log(1/(\delta\beta))}
      {\varepsilon\alpha}
 +\frac{d+\log(1/\beta)}{\alpha},
\]
\[
 R_{\mathrm{fin}}=
 \frac{\log^+|C|+\log(1/\beta)}{\varepsilon\alpha}
 +\frac{\log(1/\beta)}{\alpha},qquad
 R_{\mathrm{old}}=
 \frac{d^5\log(1/(\delta\beta))}{\varepsilon\alpha}
 +\frac{d+\log(1/\beta)}{\alpha}.
\]
The logarithmic envelope used for hidden factors is
\[
\Lambda=1+\log(e+d)+\log(e+v)+\log(e+\alpha^{-1})
 +\log(e+\beta^{-1})+\log(e+\varepsilon^{-1})
 +\log\!\bigl(e+\log(e/\delta)\bigr).
\]
The notation $N=\widetilde O(R_{\mathrm{VC}})$ means
$N\le K\Lambda^qR_{\mathrm{VC}}$ for universal $K\ge1$ and
$q\in\mathbb N_0$; no positive power of $d$, $v$, or $\log^+|C|$ may be
hidden in the tilde.
